%% kapa: exam=1 round=31 year=2020 period=2 form=A subject=appraisal_law subject_ko=감정평가관계법규 seq=38 number=38
%% source: Q-Net 감정평가사 기출문제 (cst00309) artlSeq=5207219 file=extracted/1008/5207219/A형/A형/2020년도 제31회 감정평가사 1차 시험_2교시_A형.pdf
%% conversion: kordoc 4.24.51 -> tools/gen1.py (auto); review=proofread-sonnet; flags=none
\begin{question}{38}
도시 및 주거환경정비법령상 정비구역의 해제사유에 해당하는 것은?
\begin{choices}
  \item 조합의 재건축사업의 경우, 토지등소유자가 정비구역으로 지정․고시된 날부터 1년이 되는 날까지 조합설립추진위원회의 승인을 신청하지 않은 경우
  \item 조합의 재건축사업의 경우, 토지등소유자가 정비구역으로 지정․고시된 날부터 2년이 되는 날까지 조합설립인가를 신청하지 않은 경우
  \item 조합의 재건축사업의 경우, 조합설립추진위원회가 추진위원회 승인일부터 1년이 되는 날까지 조합설립인가를 신청하지 않은 경우
  \item 토지등소유자가 재개발사업을 시행하는 경우로서 토지등소유자가 정비구역으로 지정ㆍ고시된 날부터 5년이 되는 날까지 사업시행계획인가를 신청하지 않은 경우
  \item 조합설립추진위원회가 구성된 구역에서 토지등소유자의 100분의 20이 정비구역의 해제를 요청한 경우
\end{choices}
\end{question}
\answer{4}
